
\begin{remark*}[Exposition]
The basic function chapter studies maps as relations with declared domains and
codomains. The orderings chapter adds ordered sets and order-preserving maps.
This chapter studies what happens when those two viewpoints are used together:
change the domain, change the codomain, restrict to useful pieces, and track
which properties survive.
\end{remark*}

\begin{remark*}[Changing Domains and Codomains --- Quick Reference]
\small
\begin{tabularx}{\textwidth}{l X}
\toprule
\textbf{Construction} & \textbf{Effect} \\
\midrule
Restrict domain & May improve injectivity; preserves any property checked only on surviving inputs. \\
Corestrict to image & Turns every function into a surjection onto its image. \\
Enlarge codomain & Preserves injectivity; may destroy surjectivity. \\
Extend domain & May destroy injectivity, boundedness, or monotonicity. \\
Pass to image & Gives the natural codomain for an inverse-on-image when the map is injective. \\
\bottomrule
\end{tabularx}
\end{remark*}

\subsection*{Corestriction and Same-Graph Changes}

\begin{definitionbox}{Definition (Corestriction to the Image)}
\begin{definition}[Corestriction to the Image]\label{def:corestriction-to-image}
Let $f:A\to B$ be a function. The \emph{corestriction of $f$ to its image} is
the function
\[
f_{\operatorname{im}}:A\to \im(f),
\qquad
f_{\operatorname{im}}(a)=f(a).
\]
It has the same domain and same values as $f$, but its codomain is
$\im(f)$ instead of $B$.
\end{definition}
\end{definitionbox}

\begin{remark*}[Standard quantified statement]
\[
  f_{\operatorname{im}}:A\to\im(f),
  \qquad
  \forall a\in A,\ f_{\operatorname{im}}(a)=f(a).
\]
\end{remark*}

\begin{remark*}[Interpretation]
Corestriction keeps the same rule and domain while replacing the codomain by
the image.
\end{remark*}

\begin{remark*}[Exposition]
The statement ``$f_{\operatorname{im}}$ is onto'' reads
\[
\forall y\in\im(f)\,\exists a\in A,\quad f_{\operatorname{im}}(a)=y.
\]
This is true by the definition of image: membership in $\im(f)$ already means
being equal to $f(a)$ for at least one input.
\end{remark*}

\begin{dependencies}
\begin{itemize}
  \item \hyperref[def:function]{Function}
  \item \hyperref[def:image-function]{Image of a Function}
  \item \hyperref[def:codomain]{Codomain of a Function}
\end{itemize}
\end{dependencies}

\begin{propositionbox}{Proposition (Corestriction Is Surjective)}
\begin{proposition}[Corestriction Is Surjective]\label{prop:corestriction-surjective}
For every function $f:A\to B$, the corestriction
$f_{\operatorname{im}}:A\to\im(f)$ is surjective.
\hyperref[prf:corestriction-surjective]{Go to proof.}
\end{proposition}
\end{propositionbox}

\begin{remark*}[Standard quantified statement]
\[
  \forall f:A\to B,\quad
  \operatorname{Surjective}(f_{\operatorname{im}}).
\]
\end{remark*}

\begin{remark*}[Interpretation]
Every function is onto its own image after corestriction.
\end{remark*}

\begin{dependencies}
\begin{itemize}
  \item \hyperref[def:image-function]{Image of a Function}
  \item \hyperref[def:surjective]{Surjective Function}
\end{itemize}
\end{dependencies}

\begin{figure}[htbp]
\centering
\begin{tikzpicture}[atlas, x=1cm, y=1cm, >=Stealth, every node/.style={font=\small}]
  \draw[atlasgray, line width=0.7pt] (-3.4,0) ellipse (0.7 and 1.45);
  \draw[atlasgray, line width=0.7pt] (-0.9,0) ellipse (0.95 and 1.65);
  \draw[atlasgold, line width=1pt, rounded corners=7pt] (-1.45,-0.9) rectangle (-0.35,0.95);

  \node[atlas label] at (-3.4,1.65) {$A$};
  \node[atlas label] at (-0.9,1.85) {$B$};
  \node[atlas label] at (-0.9,-1.18) {$\im(f)$};

  \node[atlas dot, label={[atlas label,left]$a_1$}] (a1) at (-3.4,0.85) {};
  \node[atlas dot, label={[atlas label,left]$a_2$}] (a2) at (-3.4,0.05) {};
  \node[atlas dot, label={[atlas label,left]$a_3$}] (a3) at (-3.4,-0.75) {};
  \node[atlas dot, label={[atlas label,right]$y_1$}] (b1) at (-0.9,0.55) {};
  \node[atlas dot, label={[atlas label,right]$y_2$}] (b2) at (-0.9,-0.55) {};
  \node[atlas dot, label={[atlas label,right]$z$}] (b3) at (-0.9,1.25) {};

  \draw[atlasblue, atlas probe, ->] (a1) -- (b1);
  \draw[atlasblue, atlas probe, ->] (a2) -- (b1);
  \draw[atlasblue, atlas probe, ->] (a3) -- (b2);

  \node[atlas label] at (1.25,0.55) {$A\to B$: not onto if $z$ is missed};
  \node[atlas label] at (1.25,-0.35) {$A\to\im(f)$: onto by definition};
\end{tikzpicture}

\caption{Changing only the codomain from $B$ to $\im(f)$ turns $f$ into a surjection.}
\label{fig:corestriction-to-image}
\end{figure}

\begin{propositionbox}{Proposition (Image Factorization)}
\begin{proposition}[Image Factorization]\label{prop:image-factorization}
Let $f:A\to B$ and let $\iota:\im(f)\hookrightarrow B$ be the inclusion map.
Then
\[
f=\iota\circ f_{\operatorname{im}}.
\]
Thus every function factors as a surjection onto its image followed by an
inclusion.
\hyperref[prf:image-factorization]{Go to proof.}
\end{proposition}
\end{propositionbox}

\begin{remark*}[Standard quantified statement]
\[
  \forall a\in A,\quad f(a)=(\iota\circ f_{\operatorname{im}})(a).
\]
\end{remark*}

\begin{remark*}[Interpretation]
Any function factors through its image by first mapping onto the image and then
including the image in the original codomain.
\end{remark*}

\begin{dependencies}
\begin{itemize}
  \item \hyperref[def:corestriction-to-image]{Corestriction to the Image}
  \item \hyperref[def:inclusion]{Inclusion Map}
  \item \hyperref[def:composition]{Composition}
\end{itemize}
\end{dependencies}

\subsection*{Restriction and Extension}

\begin{propositionbox}{Proposition (Domain Restriction Preserves Injectivity)}
\begin{proposition}[Domain Restriction Preserves Injectivity]\label{prop:restriction-preserves-injective}
If $f:A\to B$ is injective and $C\subseteq A$, then the restriction
$f|_C:C\to B$ is injective.
\hyperref[prf:restriction-preserves-injective]{Go to proof.}
\end{proposition}
\end{propositionbox}

\begin{remark*}[Standard quantified statement]
\[
  \operatorname{Injective}(f)
  \wedge C\subseteq A
  \to
  \operatorname{Injective}(f|_C).
\]
\end{remark*}

\begin{remark*}[Interpretation]
Removing inputs cannot create a collision that was not already present.
\end{remark*}

\begin{remark*}[Exposition]
The contrapositive is often the useful form: if some restriction $f|_C$ is not
injective, then $f$ is not injective. A collision inside a smaller domain is
also a collision in the larger domain.
\end{remark*}

\begin{dependencies}
\begin{itemize}
  \item \hyperref[def:restriction]{Restriction}
  \item \hyperref[def:injective]{Injective Function}
  \item \hyperref[def:subset]{Subset}
\end{itemize}
\end{dependencies}

\begin{propositionbox}{Proposition (Enlarging Codomain Preserves Injectivity)}
\begin{proposition}[Enlarging Codomain Preserves Injectivity]\label{prop:enlarge-codomain-preserves-injective}
Let $f:A\to B$ be injective and suppose $B\subseteq D$. Define
$g:A\to D$ by $g(a)=f(a)$. Then $g$ is injective.
\hyperref[prf:enlarge-codomain-preserves-injective]{Go to proof.}
\end{proposition}
\end{propositionbox}

\begin{remark*}[Standard quantified statement]
\[
  \operatorname{Injective}(f)
  \wedge B\subseteq D
  \wedge \forall a\in A,\ g(a)=f(a)
  \to
  \operatorname{Injective}(g).
\]
\end{remark*}

\begin{remark*}[Interpretation]
Injectivity depends on equality of values, not on unused elements added to the
codomain.
\end{remark*}

\begin{remark*}[Exposition]
Surjectivity is codomain-sensitive. The same rule may be onto $B$ but fail to
be onto a larger codomain $D$ because new target elements have been added.
Injectivity is different: it only asks whether two domain elements collide.
\end{remark*}

\begin{dependencies}
\begin{itemize}
  \item \hyperref[def:codomain]{Codomain of a Function}
  \item \hyperref[def:injective]{Injective Function}
  \item \hyperref[def:subset]{Subset}
\end{itemize}
\end{dependencies}
